\section{The Six Alternative Systems}
\label{sec:systems}

This section provides a brief description of each experimental system,
its core design departure from the single-member baseline, and the primary
paper where it is developed in full.

\subsection{E.1 --- Multi-Member Districts}

Multi-member districts (MMDs) replace 435 single-member districts with
approximately 145 three-member or 87 five-member districts.  Within each
district, seats are allocated proportionally to party vote share via
single transferable vote (STV) or party-list mechanisms.  A three-member
district with a 60--40 partisan split awards two seats to the majority party
and one to the minority; a five-member district with the same split awards
three to two.

The algorithmic implementation applies the same edge-weighted recursive
bisection used in the baseline \citep{dellaLibera2026dia}, but targets
$\lfloor 435/k \rfloor$ districts for member-size $k \in \{3, 5\}$.  Each
district must satisfy the population constraint that its total population
equals $k$ times the single-district quota.

MMDs are the only system in our portfolio with demonstrated empirical adoption
at scale: Australia uses three- to seven-member Senate divisions; Ireland uses
three- to five-member Dail constituencies; the UK used two-member districts
until 1950 \citep{farrell2011electoral}.  The Fair Representation Act, introduced
in multiple Congresses, would establish three- to five-member US House districts.

The primary trade-off is that larger districts reduce geographic coherence
(each representative is accountable to a more heterogeneous constituency) and
increase ballot complexity substantially.

\subsection{E.2 --- County-Based Representation}

County representation treats the 3,143 US counties and county-equivalents as
the atomic unit of representation, allocating House seats via the same
Huntington--Hill method used for state apportionment
\citep{balinski2001fair,huntington1921methods}.  Counties too small to
receive a whole seat receive fractional voting weight in Congress; counties
large enough to receive multiple seats elect multiple at-large representatives.

This system eliminates redistricting entirely: county boundaries are fixed by
state statute, not subject to manipulation by partisan legislatures.  It also
eliminates the compactness--proportionality tension within states, since seats
are allocated mechanically by population.

The principal trade-off is that fractional voting is constitutionally and
institutionally unfamiliar: the House has operated on one-member one-vote
principles since its founding.  A second objection is that county populations
vary by a factor of nearly 20,000 (Los Angeles County at 10 million versus
Loving County, Texas at approximately 64), creating extreme disparities in
the size of constituencies represented by whole-seat members.

\subsection{E.3 --- National Redistricting Without State Lines}

National redistricting treats the entire contiguous United States as a single
graph of approximately 73,000 census tracts and partitions it into 435
districts without regard to state boundaries.  Hawaii and Alaska are processed
as separate two- and one-district states respectively due to non-contiguity
with the continental geography.

The constitutional barrier is explicit: Article~I, \S2 provides that
``Representatives shall be apportioned among the several States.''  National
redistricting is therefore not a reform but a counterfactual --- a measurement
of the compactness cost imposed by federalism.  If state-based apportionment
forces districts to respect state lines, how much compactness do we sacrifice?
Paper E.3 finds approximately a 12\% improvement in mean Polsby--Popper when
state constraints are removed \citep{dellaLibera2026e3}.

The experiment also reveals that a substantial minority of current districts
cross no naturally compact boundaries: the state constraint is non-binding for
roughly 60\% of districts, which would fall entirely within one state under
national optimization anyway.  The cost is concentrated in the 40\% of
districts that currently end at awkward state borders.

\subsection{E.4 --- Partisan-Similarity Clustering}

The partisan-similarity system inverts the standard reform objective.  Instead
of creating competitive districts, it assigns high edge weights to tract pairs
with similar partisan composition, causing the METIS minimum-cut algorithm to
cluster partisan-similar tracts into the same district.  The result is a map
with a high proportion of algorithmically safe seats --- districts where the
majority-party margin is wide and stable.

This design is deliberately controversial.  The reform consensus holds that
competitive districts improve democratic accountability, forcing representatives
to be responsive to swing voters rather than to primary electorates
\citep{masket2015no}.  Partisan-similarity clustering produces the structural
opposite: near-guaranteed outcomes, primary dominance, and no electoral pressure
toward moderation.

We implement it because it demonstrates two important points.  First, extreme
partisan sorting is technically achievable through transparent algorithmic
means, not just through opaque legislative back-room dealing.  Second, if one
believes that safe seats allow legislators to focus on governance rather than
perpetual campaigning, then partisan-similarity clustering is the fair
algorithmic implementation of that preference --- and it can be contrasted with
the partisan gerrymander, which creates safe seats for one party at the expense
of the other.  Our implementation is symmetric: both parties receive their
algorithmically fair share of safe seats.

\subsection{E.5 --- Partisan-Fairness Seat Allocation}

The partisan-fairness system begins with the national two-party vote share and
constructs a map in which the number of seats won by each party exactly matches
its national vote share.  If Democrats receive 53\% of the national House vote,
they are guaranteed 230 seats; Republicans are guaranteed 205.  The algorithm
then draws compact districts subject to this seat-count constraint.

Achieving exact proportionality in a geographic system is technically demanding.
Paper E.5 shows that it requires the algorithm to accept non-compact districts
in geographically unfavorable states --- states where the natural sorting of
partisan voters makes proportionality geometrically expensive.  In practice,
the efficiency gap \citep{stephanopoulos2015partisan} provides a continuous
measure of how far a given map departs from proportionality, and E.5 optimizes
directly against that gap.

The system is constitutionally available: Congress can mandate proportionality
standards in state redistricting under Article~I, \S4, and the efficiency-gap
standard has been proposed as a statutory criterion.  The primary objection is
judicial: the Supreme Court has declined to recognize partisan symmetry as
a constitutional requirement \citep{rucho2019}, and a congressional mandate
may face separation-of-powers challenges.

\subsection{E.6 --- International Applications}

The sixth experiment applies the edge-weighted recursive bisection algorithm to
four Westminster-system countries: the United Kingdom, Canada, Australia, and
New Zealand.  These countries share the single-member district architecture
with the US but differ in boundary commission structure, redistricting frequency,
and geographic scale.

Paper E.6 finds that the algorithm transfers cleanly to each country's census
geography, producing maps with higher Polsby--Popper scores than the current
commission-drawn boundaries in all four cases.  The improvement ranges from
9\% in New Zealand (which already uses an independent commission with explicit
compactness criteria) to 21\% in Australia (where redistribution criteria
prioritize community-of-interest factors over geometric compactness)
\citep{dellaLibera2026e6}.

The international experiment serves two purposes in this synthesis.  First, it
confirms that the compactness improvements documented for the US are not
artifacts of US census geography but reflect a general property of
minimum-cut partitioning relative to commission heuristics.  Second, it
provides a comparative benchmark: the US, with no compactness criterion in
federal statute and redistricting controlled by state legislatures, starts from
a lower baseline than commission-based systems yet benefits more from
algorithmic optimization.
